% Auto-generated from report/1.md — do not edit by hand.
\section{بستر پژوهش و ضرورت پیش‌بینی محاسباتی تعاملات زیست‌مولکولی}\label{sec:1-1}


شناخت و پیش‌بینی برهم‌کنش‌های میان موجودیت‌های زیست‌مولکولی\footnote{زیست‌مولکولی (\lr{Biomolecular})} یکی از ارکان بنیادین زیست‌شناسی سامانه‌ها\footnote{زیست‌شناسی سامانه‌ها (\lr{Systems Biology})} و فرایند کشف دارو\footnote{کشف و بازکاربردپذیری دارو (\lr{Drug Discovery \& Repurposing})} به شمار می‌رود. با گسترش فناوری‌های آزمایشگاهی با توان عملیاتی بالا\footnote{غربال‌گری با توان عملیاتی بالا (\lr{High-Throughput Screening})}، حجم قابل‌توجهی از داده‌های تعاملی در پایگاه‌های داده عمومی انباشته شده است. با این وجود، فضای برهم‌کنش‌های بالقوه به‌مراتب بزرگ‌تر از فضای مشاهده‌شده است و آزمون تجربی کامل آن با محدودیت‌های هزینه و زمان مواجه است\footnote{هزینه‌های بالای مراحل پیش‌بالینی و نرخ شکست در فازهای بالینی.}. در این راستا، روش‌های محاسباتی غربال‌گری مجازی\footnote{غربال‌گری مجازی (\lr{Virtual Screening})} به‌عنوان ابزار اولویت‌بندی و کاهش فضای جست‌وجو جایگاه ویژه‌ای یافته‌اند.

تعاملات زیست‌مولکولی مورد بررسی در این پژوهش در چهار دسته ساختاری طبقه‌بندی می‌شوند:

\begin{itemize}
\item \textbf{برهم‌کنش دارو-هدف (\lr{DTI}):} پیش‌بینی پیوند فیزیکی یا مهاری میان یک مولکول دارویی و یک پروتئین هدف. این تعاملات هسته اصلی فارماکودینامیک\footnote{فارماکودینامیک (\lr{Pharmacodynamics}): مطالعه اثرات بیوشیمیایی و فیزیولوژیکی دارو بر بدن.} و کشف اهداف درمانی را تشکیل می‌دهند.
\item \textbf{برهم‌کنش دارو-دارو (\lr{DDI}):} بررسی تداخلات فارماکوکینتیکی\footnote{فارماکوکینتیک (\lr{Pharmacokinetics}): مطالعه جذب، توزیع، متابولیسم و دفع دارو.} و فارماکودینامیکی ناشی از مصرف همزمان چند دارو (پلی‌فارماسی\footnote{پلی‌فارماسی (\lr{Polypharmacy}): مصرف همزمان چندین دارو توسط یک بیمار.}).
\item \textbf{برهم‌کنش پروتئین-پروتئین (\lr{PPI}):} نگاشت شبکه فیزیکی و عملکردی پروتئین‌ها در قالب اینتراکتوم سلولی\footnote{اینتراکتوم (\lr{Interactome}): مجموعه کامل برهم‌کنش‌های مولکولی در یک سلول یا ارگانیسم.}.
\item \textbf{برهم‌کنش ژن-بیماری (\lr{GDI}):} شناسایی همبستگی‌های میان دگرگونی‌های ژنتیکی و فنوتیپ‌های بیماری‌زا بر پایه مطالعات هم‌خوانی سراسر ژنوم\footnote{مطالعات هم‌خوانی سراسر ژنوم (\lr{Genome-Wide Association Studies}).}.
\end{itemize}

نمایش یکپارچه این تعاملات در قالب شبکه‌های پیچیده زیستی، زمینه را برای به‌کارگیری الگوریتم‌های یادگیری بازنمایی روی گراف\footnote{یادگیری بازنمایی روی گراف (\lr{Graph Representation Learning}).} فراهم می‌کند.

\section{صورت‌بندی مسئله: پیش‌بینی پیوند در گراف‌های زیستی}\label{sec:1-2}


از منظر علوم رایانه، پیش‌بینی تعاملات زیست‌مولکولی به‌صورت مسئله پیش‌بینی پیوند در شبکه‌ها\footnote{پیش‌بینی پیوند (\lr{Link Prediction}).} صورت‌بندی می‌شود. شبکه زیستی به‌صورت گراف $G = (\mathcal{V}, \mathcal{E})$ مدل می‌گردد که در آن $\mathcal{V}$ مجموعه گره‌ها با کاردینالیتی $|\mathcal{V}| = N$ و $\mathcal{E} \subseteq \mathcal{V} \times \mathcal{V}$ مجموعه یال‌های مشاهده‌شده است. ماتریس مجاورت $A \in \{0, 1\}^{N \times N}$ به‌گونه‌ای تعریف می‌شود که $A_{uv} = 1$ بیانگر وجود برهم‌کنش میان گره $u$ و گره $v$ است.

هدف، یادگیری تابع نگاشت پارامتری زیر است:


\begin{equation}
\label{eq:1.1}
f_\theta: \mathcal{V} \times \mathcal{V} \to [0, 1], \qquad \hat{y}_{uv} = f_\theta(u, v \mid G)
\end{equation}


مدل‌سازی این مسئله روی گراف‌های زیستی با سه چالش ساختاری مواجه است:

\textbf{چالش اول  -  تنک بودن شدید گراف:} نسبت یال‌های شناخته‌شده به کل جفت‌های ممکن بسیار ناچیز است ($|\mathcal{E}| \ll N^2$). مدل باید در غیاب بخش عمده اتصالات، روابط پنهان را تعمیم دهد.

\textbf{چالش دوم  -  توزیع دم‌سنگین درجه گره‌ها:} تعداد اندکی گره پرتراکم (هاب) و انبوهی گره کم‌ارتباط وجود دارد \cite{barabasi1999}. این توزیع نامتقارن می‌تواند مدل‌های ساده را به سمت بهره‌برداری از محبوبیت درجه (\lr{Degree Exploitation}) سوق دهد.

\textbf{چالش سوم  -  دوگانگی ساختاری گراف‌های همگن و دوقطبی:} در شبکه‌های همگن (مانند \lr{DDI} و \lr{PPI})، تمام گره‌ها از یک نوع هستند. در شبکه‌های دوقطبی (مانند \lr{DTI} و \lr{GDI})، گره‌ها به دو دسته مجزا افراز می‌شوند ($\mathcal{V} = \mathcal{V}_1 \cup \mathcal{V}_2$) و هر یال الزاماً یک گره از $\mathcal{V}_1$ را به گره‌ای از $\mathcal{V}_2$ متصل می‌کند. در این ساختار، مسیرهای با طول زوج صرفاً گره‌های هم‌نوع را به یکدیگر وصل می‌کنند و اولین مسیر بازگشت به موجودیت مخالف، از طول فرد ($A^3$) خواهد بود. این ویژگی هندسی، پیامدهای مهمی برای طراحی عملگرهای پرش توپولوژیک دارد که در فصل دوم به تفصیل اثبات خواهد شد.

\section{مدل مرجع و محدودیت‌های شناسایی‌شده}\label{sec:1-3}


با ظهور شبکه‌های عصبی گراف\footnote{شبکه‌های عصبی گراف (\lr{Graph Neural Networks}).}، به‌ویژه شبکه‌های پیچشی گراف\footnote{شبکه‌های پیچشی گراف (\lr{Graph Convolutional Networks}).}\cite{kipf2017gcn}، پیشرفت‌های قابل‌توجهی در استخراج بازنمایی‌های برداری گره‌ها حاصل شد. مدل‌های استاندارد پیام‌ها را صرفاً در امتداد یال‌های مستقیم (۱-گام) عبور می‌دهند:


\begin{equation}
\label{eq:1.2}
H^{(l+1)} = \sigma\left(\tilde{D}^{-1/2} \tilde{A} \tilde{D}^{-1/2} H^{(l)} W^{(l)}\right)
\end{equation}


در سال ۲۰۲۰، هوانگ و همکاران مدلی تحت عنوان \lr{SkipGNN} معرفی کردند \cite{huang2020skipgnn}. ایده محوری این مدل، تزریق «شباهت پرشی» از طریق ساخت گراف پرش ۲-گام بود:


\begin{equation}
\label{eq:1.3}
A_s = \mathrm{sign}(A A^T)
\end{equation}


\lr{SkipGNN} پیام‌ها را به‌صورت همزمان روی گراف اصلی $A$ و گراف پرش $A_s$ عبور داده و با دیکودر خطی جفت‌ها را امتیازدهی می‌کرد. این مدل با گزارش \lr{AUPRC} به میزان $0.928$ روی دیتاست \lr{DTI}، عملکرد قابل‌توجهی نسبت به روش‌های پیشین گزارش کرد.

بازبینی دقیق تئوریک و تجربی در این پژوهش، چهار محدودیت ساختاری و روش‌شناختی در \lr{SkipGNN} شناسایی کرد:

\textbf{محدودیت اول  -  بن‌بست هندسی در گراف‌های دوقطبی:} در گراف‌های دوقطبی \lr{DTI} و \lr{GDI}، ماتریس $A^2$ ساختاری بلوکی-قطری دارد. کانال پرش \lr{SkipGNN} پیام‌ها را فقط بین گره‌های هم‌نوع (دارو-دارو یا ژن-ژن) ردوبدل می‌کند و اطلاعاتی درباره تارگت‌های کاندید از نوع مخالف انتقال نمی‌دهد. اثبات کامل این موضوع در بخش ۲-۴-۱ ارائه می‌شود.

\textbf{محدودیت دوم  -  حذف اطلاعات چگالی بر اثر باینری‌سازی:} استفاده از تابع $\mathrm{sign}$ در ساخت گراف پرش، تمایز میان جفت‌هایی با تعداد متفاوت همسایه مشترک را از بین می‌برد. جفتی با ۵۰ همسایه مشترک و جفتی با ۱ همسایه مشترک، وزن یکسانی ($1.0$) دریافت می‌کنند.

\textbf{محدودیت سوم  -  تنگنای ظرفیت دیکودر:} دیکودر \lr{SkipGNN} صرفاً یک لایه خطی روی الحاق بردارها $[E_u \parallel E_v]$ است. این دیکودر فاقد مولفه‌های تعاملی ضربی و فاصله‌ای لازم برای بازنمایی برهم‌کنش‌های بیوشیمیایی متقارن است.

\textbf{محدودیت چهارم  -  سوگیری ساختاری در پروتکل ارزیابی:} پروتکل ارزیابی مقاله مرجع، جفت‌های منفی را به‌صورت تصادفی و یکنواخت انتخاب می‌کرد. با توجه به توزیع دم‌سنگین درجه در گراف‌های زیستی، بخش عمده منفی‌های تصادفی از گره‌های کم‌درجه تشکیل می‌شوند. این امر تمایز یال‌های واقعی از منفی‌ها را به‌طور مصنوعی ساده می‌کند و ارزیابی قدرت تعمیم ساختاری مدل را دشوار می‌سازد. شواهد تجربی این پژوهش (فصل پنجم) نشان می‌دهد که عملکرد \lr{SkipGNN} در آزمون یکنواخت ($0.929$ در \lr{DTI}) به‌طور قابل‌توجهی بالاتر از عملکرد آن در آزمون با منفی‌های هم‌درجه ($0.695$) است. این الگو با یافته‌های لی و همکاران \cite{li2023evaluating} درباره محدودیت بنچمارک‌های مبتنی بر منفی‌های تصادفی هم‌راستا است.

\section{پرسش‌ها و فرضیات پژوهش}\label{sec:1-4}


این پروژه حول پنج پرسش پژوهشی و فرضیات متناظر شکل گرفته است:

\textbf{پرسش اول (\lr{RQ1})  -  چگالی توپولوژیک پرش:} آیا باینری‌سازی گراف پرش توسط عملگر $\mathrm{sign}$ موجب حذف اطلاعات شدت تعاملات می‌شود؟

\emph{فرضیه ۱ (\lr{H1}):} جایگزینی ماتریس باینری با ماتریس پیوسته تخصیص منابع \cite{zhou2009predicting} به فرم $W = A D^{-1} A^T$، تمایز میان جفت‌های دارای مسیرهای غنی و ضعیف را حفظ کرده و از اشباع گرادیان در گره‌های هاب جلوگیری می‌کند.

\textbf{پرسش دوم (\lr{RQ2})  -  دوگانگی ساختار دوقطبی:} چگونه می‌توان در گراف‌های دوقطبی، پل اطلاعاتی میان موجودیت‌های ناهم‌نوع برقرار کرد؟

\emph{فرضیه ۲ (\lr{H2}):} تزریق صریح مسیرهای بازگشت ۳-گام به فرم $W_3 = W_2 D^{-1} A$ به‌عنوان کانال پیام‌رسانی، بن‌بست هندسی گام‌های زوج را جبران می‌نماید.

\textbf{پرسش سوم (\lr{RQ3})  -  ظرفیت دیکودر:} آیا دیکودرهای خطی مبتنی بر الحاق ساده در تفکیک منفی‌های سخت ناتوانند؟

\emph{فرضیه ۳ (\lr{H3}):} استفاده از دیکودر تانسوری تعاملی ۴‌گانه متشکل از الحاق، ضرب هادامارد و تفاضل قدرمطلقی $[E_u \parallel E_v \parallel E_u \odot E_v \parallel |E_u - E_v|]$ همراه با نرمال‌سازی دسته‌ای\footnote{نرمال‌سازی دسته‌ای (\lr{Batch Normalization}).}\cite{ioffe2015batch}، تقارن و غیرخطی بودن تعاملات را بازنمایی کرده و دقت تفکیک را بهبود می‌بخشد.

\textbf{پرسش چهارم (\lr{RQ4})  -  تنظیم‌سازی فضا و پایداری بازنمایی:} چگونه می‌توان از فروپاشی بازنمایی\footnote{فروپاشی بازنمایی (\lr{Representation Collapse}): پدیده‌ای که در آن امبدینگ‌های گره‌های مختلف به یک بردار یکسان همگرا می‌شوند.} و بیش‌برازش بر گره‌های هاب جلوگیری کرد؟

\emph{فرضیه ۴ (\lr{H4}):} اعمال قید خودنظارتی مقابله‌ای متقارن (\lr{InfoNCE}) میان نمای گراف اصلی و نمای گراف پرش، خاصیت یکنواختی در فضای بازنمایی\footnote{یکنواختی در فضای بازنمایی (\lr{Uniformity}).}\cite{wang2020uniformity} را ارتقا می‌دهد.

\textbf{پرسش پنجم (\lr{RQ5})  -  اعتبار پروتکل ارزیابی:} آیا عملکرد گزارش‌شده در معیارهای یکنواخت، بازتاب‌دهنده قدرت تعمیم ساختاری مدل است؟

\emph{فرضیه ۵ (\lr{H5}):} ارزیابی مدل‌ها روی بانک آزمون با منفی‌های هم‌درجه، سوگیری محبوبیت گره‌ها را حذف کرده و معیاری پایاتر برای سنجش تعمیم ساختاری ارائه می‌دهد.

\section{مدل‌های پیشنهادی و نوآوری‌های پژوهش}\label{sec:1-5}


برای پاسخ به پرسش‌های فوق، چهار مدل در چارچوب یک پایپ‌لاین یکپارچه طراحی و پیاده‌سازی شدند. هر یک از این مدل‌ها بر بُعد مشخصی از محدودیت‌های شناسایی‌شده تمرکز دارد:

\textbf{مدل ۱  -  \lr{AMS-SkipGNN} (\lr{Attentive Multi-Scale SkipGNN}):} مدل اصلی پژوهش که سه محدودیت اول را به‌طور همزمان هدف قرار می‌دهد. این مدل پرش باینری را با پرش پیوسته تخصیص منابع جایگزین می‌کند، جمع ساده لایه‌ای را با واحدهای گیتینگ تطبیقی باقی‌مانده تعویض می‌نماید و دیکودر خطی را با دیکودر تانسوری تعاملی ۴‌گانه جایگزین می‌سازد.

\textbf{مدل ۲  -  \lr{SkipGATv2}:} این مدل محدودیت انتشار ثابت لاپلاسی را هدف قرار می‌دهد. با جایگزینی کانوولوشن ثابت با مکانیزم توجه پویا لبه‌ای \cite{brody2022gatv2}، شبکه قادر است یال‌های کم‌اطلاع را به‌صورت خودکار کم‌اهمیت کند. این مدل در گراف همگن \lr{PPI} عملکرد قابل‌توجهی ثبت کرده و در بخش ۵-۲ تحلیل می‌شود.

\textbf{مدل ۳  -  \lr{ThreeHopSkipGNN}:} این مدل به‌طور اختصاصی بن‌بست هندسی گراف‌های دوقطبی (محدودیت اول) را هدف قرار می‌دهد. با ادغام موازی سه کانال ۱-گام، ۲-گام هم‌نوع و ۳-گام دوقطبی و ترکیب آن‌ها از طریق گیت سافت‌مکس یادگیرنده، دسترسی مستقیم انکودر به اهداف کاندید در گراف‌های دوقطبی فراهم می‌شود.

\textbf{مدل ۴  -  \lr{ContrastiveSkipGNN}:} این مدل محدودیت چهارم (پایداری فضای بازنمایی) را هدف قرار می‌دهد. با ترکیب یادگیری بانظارت پیش‌بینی پیوند و قید خودنظارتی مقابله‌ای متقارن، از انحراف و چسبندگی امبدینگ‌های هاب جلوگیری می‌شود.

\textbf{نکته شفاف‌سازی:} مدل \lr{AMS} عمیق‌ترین اعتبارسنجی ساختاری (نردبان ابلیشن کامل، بخش ۴-۷) را در میان چهار مدل دارد و در دیتاست \lr{DDI} به‌وضوح برترین مدل عصبی است. با این حال، از نظر امتیاز خام \lr{Hard AUPRC}، در سه دیتاست دیگر (\lr{DTI}، \lr{PPI}، \lr{GDI}) با \lr{SkipGATv2} هم‌سطح یا در محدوده انحراف معیار آن قرار دارد. مقایسه کامل و دقیق آماری در فصل پنجم و بحث در فصل ششم ارائه می‌شود.

\subsection{نوآوری‌های روش‌شناختی}\label{sec:1-5-1}


علاوه بر نوآوری‌های معماری، دو کمک روش‌شناختی در این پژوهش ارائه شده است:

\textbf{پروتکل ارزیابی دوگانه (\lr{Dual-Bank Protocol}):} ارزیابی همزمان بر روی بانک یکنواخت (برای بازتولید شرایط مقایسه با مقاله مرجع) و بانک سخت آگاه از نوع موجودیت با افراز چارک‌های درجه‌ای.

\textbf{پایپ‌لاین بدون نشت ترنسداکتیو:} جداسازی کامل ماتریس مجاورت آموزشی از یال‌های اعتبارسنجی و آزمون، انتخاب آستانه بهینه ($\tau^*$) صرفاً روی داده‌های اعتبارسنجی و انجماد آن روی آزمون\footnote{جلوگیری از نشت اطلاعات در انتخاب آستانه (\lr{Threshold Leakage Prevention}).}، و اصلاح ناسازگاری‌های فرمت تنک/متراکم در کد مرجع.

\subsection{انتخاب آگاهانه ویژگی‌های ورودی}\label{sec:1-5-2}


در این پژوهش، ویژگی‌های ورودی گره‌ها به بردارهای هویت تنک (\lr{One-Hot}) محدود شده‌اند. این انتخاب یک تصمیم روش‌شناختی آگاهانه است: هدف اصلی، سنجش ظرفیت خالص معماری‌های پیام‌رسانی در کشف الگوهای توپولوژیک پیوند است، نه ارزیابی قدرت مدل‌های زبانی پیش‌آموزش‌دیده در پردازش ویژگی‌های بیوشیمیایی. افزودن ویژگی‌های غنی (مانند اثرانگشت مورگان یا امبدینگ‌های \lr{ESM-2}) می‌تواند اثر معماری را با اثر ویژگی‌ها مخلوط کند و تفکیک سهم هر یک را دشوار سازد. این محدودیت به‌صراحت در فصل ششم به‌عنوان مسیر کار آینده ذکر شده است.

\section{دستاوردهای عددی شاخص}\label{sec:1-6}


نتایج تجربی کامل در فصل پنجم ارائه می‌شود. در اینجا، خلاصه‌ای از بهبودهای کلیدی نسبت به مدل مرجع \lr{SkipGNN} در آزمون سخت (\lr{Hard AUPRC}) گزارش می‌شود. مقادیر به فرم $\text{\lr{Mean}} \pm \text{\lr{SD}}$ در میان ۳ سید هستند. بهبود مطلق به واحد درصد (\lr{percentage points}) بیان شده است:


\begin{table}[htbp]
\centering
\footnotesize
\caption{بهبود مطلق نسبت به \lr{SkipGNN} در \lr{Hard AUPRC} (واحد درصد)}
\label{tab:1-1}
\begin{tabular}{lllll}
\toprule
دیتاست & \lr{SkipGNN} & مدل‌های پیشنهادی برتر & بهبود مطلق & وضعیت آماری \\
\midrule
\lr{DDI} & $0.724 \pm 0.005$ & \lr{AMS}: $0.912 \pm 0.003$ & $+0.188$ & برتری بدون هم‌پوشانی بازه‌ها \\
\lr{PPI} & $0.622 \pm 0.012$ & \lr{SkipGATv2}: $0.778 \pm 0.006$ & $+0.156$ & برتری بدون هم‌پوشانی بازه‌ها \\
\lr{GDI} & $0.720 \pm 0.012$ & \lr{SkipGATv2}: $0.837 \pm 0.003$ & $+0.117$ & برتری مرزی (بازه‌ها به‌سختی جدا) \\
\lr{DTI} & $0.695 \pm 0.012$ & \lr{AMS}، \lr{SkipGATv2} و \lr{3-hop} هم‌سطح & $+0.10$ تا $+0.11$ & هم‌سطح آماری (بازه‌های هم‌پوشان) \\
\bottomrule
\end{tabular}
\end{table}


\textbf{توضیح:} در دیتاست \lr{DTI}، سه مدل پیشنهادی  -  \lr{AMS} ($0.798 \pm 0.011$)، \lr{SkipGATv2} ($0.804 \pm 0.005$) و \lr{3-hop} ($0.803 \pm 0.015$)  -  بازه‌های انحراف معیار هم‌پوشانی دارند و طبق معیار عدم هم‌پوشانی بازه‌ها (بخش ۵-۱-۵)، از نظر آماری قابل تفکیک نیستند.

شایان ذکر است که در دیتاست \lr{GDI}، هوریستیک تحلیلی تخصیص منابع (بدون پارامتر یادگیرنده) به امتیاز $0.842 \pm 0.001$ دست یافته که بالاتر از تمام مدل‌های عصبی است. این یافته به‌عنوان یک بینش مهم درباره محدودیت‌های ذاتی مدل‌های عصبی در گراف‌های بسیار بزرگ و تنک با ویژگی‌های هویتی، در فصل ششم به تفصیل بحث می‌شود.

\section{ساختار فصول گزارش}\label{sec:1-7}


ادامه گزارش در پنج فصل به شرح زیر سازمان‌دهی شده است:

\textbf{فصل دوم} به مبانی نظری شبکه‌های عصبی گراف، پیش‌بینی پیوند و نقد ساختاری مدل مرجع \lr{SkipGNN} اختصاص دارد. اثبات ریاضی بن‌بست بلوکی ماتریس $A^2$ در گراف‌های دوقطبی در این فصل ارائه می‌شود.

\textbf{فصل سوم} مشخصات چهار دیتاست، پایپ‌لاین بدون نشت داده، فرمولاسیون عملگرهای پرش توپولوژیک و پروتکل ارزیابی دوگانه را تشریح می‌کند.

\textbf{فصل چهارم} فرمولاسیون ریاضی دقیق چهار مدل پیشنهادی (\lr{AMS}، \lr{SkipGATv2}، \lr{ThreeHop}، \lr{Contrastive}) و خطوط مبنا را ارائه می‌دهد.

\textbf{فصل پنجم} شامل پروتکل آموزش، ماتریس کامل نتایج عددی، تحلیل ده نمودار و بررسی فضاهای نهفته با \lr{t-SNE} است.

\textbf{فصل ششم} به بحث، تبیین چرایی نتایج، تحلیل محدودیت‌ها و موارد شکست، پاسخ به پرسش‌های پژوهش و ترسیم مسیرهای آینده می‌پردازد.
